\section{Release Behavior as Merge Evidence}

Release is where evidence debt becomes user-visible. A pull request can pass tests, look well reviewed, and still be unsafe to ship if the repository cannot identify the version, commit, tag, artifact, approval, security evidence, and recovery path that define the release. Jankurai therefore treats release governance as part of the merge-witness problem rather than as a separate checklist. HL5 does not ask a team to use one release tool. It asks the repository to prove that shipped artifacts are bound to source and can be understood, verified, and recovered.

The release evidence boundary is deliberately concrete. A release-capable repository should expose a version source, changelog, release process document, release automation or command policy, artifact integrity policy, and rollback guidance. The version can be SemVer, CalVer, a package version, a container tag, or a deployment ID, but it must be unique and ordered enough for reviewers and users to reason about compatibility. SemVer's immutability rule is the right default for public packages: once released, the content of that version is not modified; a change gets a new version \cite{semverSpec2026}. Supply-chain standards push the same direction from the artifact side: checksums, SBOMs, provenance, signatures, and attestations make artifacts inspectable instead of merely trusted by convention \cite{slsaSpec2026,spdxSpecifications2026}.

The inexcusable behaviors are the ones that destroy that binding. Retagging a release, force-pushing or deleting release refs, overwriting published assets, shipping from a dirty local tree, skipping tests or package verification, publishing only a mutable \texttt{latest} tag, hiding breaking changes from the changelog, using broad long-lived publishing tokens, running untrusted pull-request code in privileged release jobs, packaging \texttt{.env} or local credential files, and shipping without rollback guidance all convert a release into an unreviewable assertion. These failures are not matters of team taste. They remove the evidence that would let a later human, agent, incident responder, or downstream consumer reconstruct what happened.

Jankurai maps that into two release controls. \texttt{HLT-025-RELEASE-READINESS-GAP} fails repositories that claim release or publish capability without the expected release structure and launch-gate evidence: security, backup, monitoring, rollback, and abuse-control receipts. \texttt{HLT-037-RELEASE-BAD-BEHAVIOR} is detector-backed and rejects high-confidence executable hazards in release scripts, package scripts, and workflows: mutable tags and assets, skipped verification, mutable latest-only publication, secret-bearing artifacts, unverified GitHub release creation, missing checksum/SBOM/provenance/signature/attestation evidence, and privileged publishing from untrusted triggers.

This is not a claim that Jankurai proves the released software is semantically correct. It is a narrower release-evidence claim. A merge witness can bind changed paths to owners and proof lanes; a release witness must additionally bind the approved commit to version identity, artifact identity, supply-chain evidence, publication command, monitoring, and recovery. Existing CI, package managers, Git hosting platforms, SLSA tooling, SBOM generators, signing systems, and deployment platforms already emit many of those signals. Jankurai's role is to require that the signals exist, remain current, and participate in the merge or release decision instead of being reconstructed after an incident.
